HD 181068 was one of the brightest targets which was continuously observed
during the almost 4-year-long primary mission of NASA's exoplanet-hunter
Kepler space telescope. The spacecraft observed ${\approx}3-4 \times 10^{-3}$
magnitude dimmings every 0.453 days. (Note: the even dimmings were slightly
smaller amplitude than the odd ones.) Furthermore, additional 0.007
magnitude, 2.3-day-long dimmings were detected every 22.7 days.

The correct explanation of this very unusual photometric behaviour was given
by Hungarian astronomers. They found that HD 181068 is a compact hierarchical
triple stellar system seen almost edge-on.

Hierarchical triple star systems consist of three stars; A, B, and C. Two of
these stars (B and C) form an inner or close stellar binary system, whilst
the outer component (star A) orbits at a distance from the inner system
significantly larger (usually orders of magnitude) than the semi-major axis
of the inner system. The schematic view of an example of a triple star system
is illustrated in Fig.~2.1 (the schematic pole-on view of a hypothetical
hierarchical triple stellar system; the black arrow is directed towards the
Earth; the thick segments of the three orbits represent the stars' orbital
arcs during an outer eclipse).

Mathematically, the motion of a hierarchical triple system can be well
approximated with two unperturbed Keplerian two-body motions; (1) Keplerian
motion of the inner binary. (2) The centre of mass of this close binary and
the third star revolves on a second Keplerian orbit, ``outer binary''.

In this problem, stars B and C form a $P_1 = 0.9056768$-day-period eclipsing
binary, while the centre of mass of these stars with star A forms the
$P_2 = 45.4711$-day-period outer binary. As the orbital plane of this outer
orbit is seen almost edge-on from the Kepler spacecraft (and from the Earth),
during their revolution on the outer orbit, stars B and C eclipse not only
each other, but also star A or, a half outer revolution later are eclipsed by
it, causing the extra dimmings.

\begin{description}
\item[i.] \emph{Determination of the physical stellar sizes (and other
  quantities) from the geometry of the eclipses}

  These assumptions are used throughout this section: (1) both the inner and
  outer orbits are exactly circular, (2) the orbital planes of the inner and
  outer orbits are identical, and (3) this plane is seen exactly edge-on
  (i.e.\ $i_1 = i_2 = 90^\circ$ and $i_{\mathrm{rel}} = 0^\circ$). Let's
  consider the extra dimmings, which are central eclipses (i.e.\ either
  occultations or transits -- annular eclipses), therefore, these events
  have four contacts. In the case of an ordinary eclipsing binary (or of a
  transiting exoplanet) at the times of the outer contacts the sky-projected
  disks of the two objects connect with each other at one point from
  outside, while at the inner contacts they approach each other from inside.
  While this last statement is also valid for outer eclipses, the situation
  becomes more complex, because instead of two, three stars are involved
  into the eclipses. However, despite this fact, we can certainly define the
  times of each contact from the light curve, and furthermore, we can also
  decide explicitly which member of the inner binary is involved in a given
  contact. (The other member, of course, is always star A.)

  In the table below, the accurate times of some contacts of different
  eclipses observed by the Kepler spacecraft, the contact types and the
  stars are documented. Time is expressed in barycentric Julian Days (BJD).

  \begin{center}
  \begin{tabular}{cccccc}
  event no. & contact & stars & BJD & $\varphi_1$ & $\varphi_2$ \\
  1 & I & A, B & 2455476.1096 & & \\
    & II & A, C & 2455476.4245 & & \\
    & III & A, B & 2455477.9677 & & \\
    & IV & A, B & 2455478.4722 & & \\
  2 & I & A, B & 2455521.5217 & & \\
  3 & III & A, C & 2455568.9434 & & \\
  4 & I & A, C & 2455612.4733 & & \\
    & III & A, C & 2455614.3571 & & \\
  5 & III & A, B & 2455659.9241 & & \\
    & IV & A, C & 2455660.2422 & & \\
  \end{tabular}
  \end{center}

  \begin{description}
  \item[(a)] Given that $T_{01} = 2455051.2361$ and
    $T_{02} = 2455522.7318$ denote the time of an inferior conjunction of
    the inner and outer binaries respectively (i.e.\ that time, when, from
    the perspective of the observer, star C eclipses star B, and when star A
    eclipses the centre of mass of stars B and C.)

    Define
    \[ \varphi_1(t) = \{(t - T_{01})/P_1\}, \quad \mbox{and} \quad
       \varphi_2(t) = \{(t - T_{02})/P_2\} \]
    as the photometric phases of the inner and outer binaries respectively.
    $\{x\}$ denotes the decimal part of the real number $x$. If
    $\{x\} < 0$, use $\{x\} + 1$ instead. Calculate the phases for the times
    of the tabulated contact events and write the answers in the appropriate
    columns of the table on the answer sheet. Round your answers to four
    decimal places. \hfill[10\,p]
  \item[(b)] Determine, whether star A, or the close binary (i.e.\ stars B
    and C) were closer to the observer during each eclipsing event. Write
    your answer in the table on the answer sheet. \hfill[5\,p]
  \item[(c)] Using the table above, calculate (1) the dimensionless radius
    of each star relative to the semi-major axis of the outer orbit
    ($R_{\mathrm{A,B,C}}/a_2$), (2) the ratio of the semi-major axes of the
    two orbits ($a_1/a_2$) and (3) the mass ratio of stars B and C
    ($q_1 = m_{\mathrm{C}}/m_{\mathrm{B}}$). Hint: use at least four decimal
    place accuracy in your calculations. Be cautious, it may not be possible
    to use all theoretically possible contact combinations with a given
    limited accuracy of time data. \hfill[30\,p]
  \item[(d)] Based on the results obtained above, calculate the outer mass
    ratio ($q_2 = m_{\mathrm{BC}}/m_{\mathrm{A}}$). \hfill[8\,p]
  \end{description}

\item[ii.] \emph{Dynamical determination of the stellar masses using radial
  velocity (RV) and eclipse timing variations (ETV) measurements}

  To obtain RV data, ground-based spectroscopic follow up observations were
  carried out with four different instruments. Only the lines of star A were
  detectable in all spectra. Plotting all the measurements against time, the
  RV curve was nicely fitted in the following form:
  \[ V_{\mathrm{rad,A}} = V_\gamma + K_{\mathrm{A}}\sin\phi_{\mathrm{RV}}, \]
  where $V_\gamma$ is the systemic velocity and $K_{\mathrm{A}}$ is the
  velocity amplitude:
  \[ V_\gamma = 6.993 \pm 0.011\,\mathrm{km\,s^{-1}}, \quad
     K_{\mathrm{A}} = 37.195 \pm 0.053\,\mathrm{km\,s^{-1}}, \]
  \[ P_2 = 45.4711 \pm 0.0002\,\mathrm{d}, \quad
     \phi_{\mathrm{RV}} = \frac{2\pi}{P_2}
     \left[t - (2455522.7318 \pm 0.0095)\right]. \]
  Furthermore, the researchers determined the mid-times of the regular
  eclipses of the close binary (formed by stars B and C), and found that the
  occurrence of for example the eclipsing minima belonging to the $N$th
  orbital revolution can be described by the simple expression:
  \[ T_N = T_0 + P_1 N + A_{\mathrm{ETV}}
     \sin\left(\frac{2\pi}{P_2}P_1 N + \phi_0\right), \]
  where
  \[ T_0 = \mathrm{BJD}\ 2455051.23607 \pm 5\times10^{-5}, \quad
     P_1 = 0.9056768 \pm 3\times10^{-7}\,\mathrm{d}, \]
  \[ A_{\mathrm{ETV}} = 0.001446 \pm 0.000110\,\mathrm{d}, \quad
     \phi_0 = -0.76779 \pm 0.01937\,\mathrm{rad}. \]
  In this expression $A_{\mathrm{ETV}}$ is the amplitude of the eclipse
  timing variation, $T_0$ denotes the mid-eclipse time of the reference
  (zeroth) primary eclipse, and $N$ is the cycle number, which is an integer
  for primary eclipses (i.e.\ when the slightly fainter star C eclipses star
  B), and half-integer for secondary ones (i.e.\ when star B eclipses star
  C).

  Determine (1) again the mass ratio
  ($q_2 = m_{\mathrm{BC}}/m_{\mathrm{A}}$) of the centre of mass of the
  inner binary and star A using only the results obtained in point ii., (2)
  the mass of component A ($m_{\mathrm{A}}$) and (3) the total mass of the
  inner, close binary ($m_{\mathrm{BC}}$). Calculate the errors for (1),
  (2), and (3) in masses. Hint: you can save much time by expressing the
  masses in solar mass and the orbital separations either in solar radius or
  au. \hfill[22\,p]

\item[iii.] Using results obtained in questions 1 and 2, determine the
  masses of stars B and C respectively and calculate the physical dimensions
  of all three stars (i.e.\ stellar radii in physical units). \hfill[15\,p]
\end{description}
